\section{Quadratic error estimate}\label{sec:error}
Define the exact error $\err_n=x_n-\rtwo$. By the theorem it is positive, and
\eqref{eq:error-identity} yields the exact identity
\begin{equation}\label{eq:quadratic}
\err_{n+1}=\frac{\err_n^2}{2x_n}.
\end{equation}
Since $\rtwo<x_n\leq2$, this implies
\[\frac14\err_n^2\leq\err_{n+1}\leq\frac{1}{2\rtwo}\err_n^2.\]
In particular,
\[\lim_{n\to\infty}\frac{\err_{n+1}}{\err_n^2}=\frac{1}{2\rtwo},\]
which is the usual meaning of quadratic convergence with a positive asymptotic
error constant. This ratio statement concerns exact errors, all of which are
nonzero; it is not a claim about errors rounded to zero by a computer.

A quantitative bound follows by setting $q_n=\err_n/(2\rtwo)$.
Then $q_{n+1}\leq q_n^2$, and induction gives
\[\err_n\leq2\rtwo\left(\frac{2-\rtwo}{2\rtwo}\right)^{2^n}.\]
The base is less than one. The double exponent explains why only a few steps
are needed to reach ordinary machine precision. Once rounding errors are of
the same size as the exact truncation error, the recurrence for exact errors
no longer models the computed trajectory accurately. This distinction is
essential when interpreting a plot with a logarithmic axis.
